% Generated draft: requires mathematical and source review.
\documentclass[11pt]{article}
\usepackage[margin=25mm]{geometry}
\usepackage{fontspec}
\setmainfont{Latin Modern Roman}
\usepackage{amsmath,amssymb,mathtools,unicode-math}
\setmathfont{Latin Modern Math}
\usepackage{enumitem,needspace,xurl,hyperref}
\setlength{\parindent}{0pt}
\setlength{\parskip}{5pt}
\setlength{\emergencystretch}{3em}
\newcommand{\N}{\mathbb{N}}
\newcommand{\Z}{\mathbb{Z}}
\newcommand{\Q}{\mathbb{Q}}
\newcommand{\R}{\mathbb{R}}
\newcommand{\C}{\mathbb{C}}
\newcommand{\F}{\mathbb{F}}
\newcommand{\A}{\mathbb{A}}
\newcommand{\PP}{\mathbb{P}}
\newcommand{\E}{\mathbb{E}}
\newcommand{\eps}{\varepsilon}
\newcommand{\dd}{\,\mathrm{d}}
\begin{document}
\section*{UnsolvedMath ID: 3100011}
\textbf{Source problem code:} catalogue entry
\subsection{Definitions and mathematical statement}
Straus asks whether every polygonal room in the plane is illuminable from some point. The source supplies no definition of illumination, and its ``See this'' link target is missing. This omission changes the mathematical question.

For direct visibility in a polygonal region \(P\), an interior source \(p\) illuminates \(q\) when \([p,q]\subseteq P\). The one-source assertion is then
\[
\forall P\quad\exists p\in\operatorname{int}P\quad
\forall q\in\operatorname{int}P:\ [p,q]\subseteq P,
\]
which is the nonempty-kernel/star-shapedness property. This is not the same as illumination by reflected light.

In a specular billiard interpretation, \(q\) is illuminated from \(p\) if there is a path consisting of finitely many straight segments from \(p\) to \(q\), remaining in the room and reflecting at relative interiors of wall segments with equal incidence and reflection angles; vertex collisions are normally excluded. The corresponding target would be
\[
\forall P\quad\exists p\in\operatorname{int}P\quad
\forall q\in\operatorname{int}P:\quad
\text{there exists such a reflected path from }p\text{ to }q.
\]
The number of reflections may depend on \(q\). This is explicit conditional notation, not a claim that the source fixes this model. Allowing holes, boundary target points, corner reflections, or an almost-everywhere rather than every-point conclusion also changes the target.

The supplied OPEN-TRIAGE summary conflates these interpretations. It states both that every simple polygon is illuminable and that non-star-shaped polygons give negative one-source examples, invokes the finite-guard art-gallery theorem, and ends in a truncated ``Toben--penrose'' discussion. These statements do not supply a coherent resolution of the source's illumination problem. No single-point reflected-light conclusion is inferred from direct visibility or finite-guard results.
\subsection{Short English statement}
Does every polygonal room admit one point that illuminates the entire room? The source does not specify direct or reflected light, and its contradictory summary cannot settle the intended question.
\subsection{Sources}
\begin{itemize}
\item Local UnsolvedMath registry, permanent ID \(3100011\), exact source code \texttt{catalogue entry}, ``Is every polygonal room in the plane illuminable from some point,'' attributed to Straus.
\item Cooper, \emph{Combinatorial Problems I Like}, October 2020 snapshot, Geometry, source-order bullet 11; \url{https://people.math.sc.edu/cooper/combprob.html}. HTML extraction; recorded access 2026-07-29. The additional ``See this'' target is not supplied.
\item The research summary mentions star-shapedness, the art-gallery theorem, and a truncated ``Toben--penrose'' item without complete authoritative citations or a coherent resolution of the stated problem.
\end{itemize}
\end{document}
